\documentclass[a4]{article}
\usepackage{gnuplottex}
\usepackage{csvsimple}
\usepackage{subcaption}
\usepackage{amsmath}

\newif\ifc
\cfalse
\newif\ifjava
\javafalse
\newif\ifpython
\pythontrue

\title{COMP26120 Lab 3 Report}
\author{Louise A. Dennis}
\begin{document}
\maketitle

\section{A Note on Testing}

We tested your hash set implementation on all the simple tests and the collision\_tests tests.  We did this both with the default hash set size (509 unless you had changed it), with an initial size of 11 excluding simple 7 \& 8 (chosen to avoid rehashing in most cases but with a high likelihood of collisions) and with an initial size of 3 (to force rehashing).  Most people's implementations passed the default size, some had failures, usually related to collision resolution with the size of 11, quite a lot of solutions failed with an initial size of 3 - in particular if the load factor was set to something like 75\% then for the examples in simple with only three items in the dictionary, rehashing was not called but the dictionary was full when find was called and in some cases find then failed to terminate since there were no empty cells in the dictionary.

For basic testing we expected people's implementations to pass both default size and initial size 11.  For rehashing testing we also expected them to pass for size 3.

For binary trees we ran your implementations on the simple tests.

We encountered quite a lot of errors in the implementation of print\_set and print\_stats, particularly for binary trees where these were sometimes called on null elements (i.e., left or right branches of the tree that had not been created).  Where possible we fixed these before running the binary trees tests.  

\section{Part A: Short Answers to Quiz}

\begin{quote}
What did you implement as your hash function (give a formula to describe what it computes)? How “good” is it (justify your answer)? (No more than 200 words).  

If you have implemented more than one hash function then this should describe the hash function used in mode 0 (and the one that is used first, if you are using double hashing as your collision resolution strategy). 
\end{quote}

Given a string $k$ of length $l$ where the ASCII value of character $i$ in the string is  $k[i]$ and dictionary size $N$.  This computes:

$$k[0] \times 41 k[1] \times 41^2 k[2] \times \ldots 41^{l-1} k[l - 1] \text{ mod } N$$

This hash function will produce a nice spread of values -- the position of letters in the string is important to the value produced and the use of powers of 41 means a good range of numbers will be used even for quite short strings.  The use of powers of 41 related to the position of the letter in the string mean it avoids hashing symmetric strings (i.e., anagrams) to the same value.

{\bf Marking Feedback:}  Where most people lost marks here was reporting that their hash function was good (or bad) but failing to explain why -- i.e., not discussing the spread of values and its handling of symmetry.  Some people lost marks because they used a built-in hash function - e.g., \texttt{hashCode} in Java and couldn't explain how it worked - in this case simply saying it was a polynomial hashing function was insufficient since it wasn't clear you knew what that meant.

\begin{quote}
What is your collision resolution strategy? What are the advantages and disadvantages of this strategy (this should include any factors that should be considered in the implementation for it to be correct (e.g., limitations on the choices for hash table size or hash function))?  How do find and insert for hash sets work using this strategy?  How have you implemented them?  Illustrate your explanation with samples of your code - this may be tided up a bit for presentation but should clearly be the same as the code you have submitted and should cover the the operation of collision resolution, and how it is used by both find and insert (No more than 300 words, not counting the code snippets).
\end{quote}

I have implemented linear probing.  If my choice of cell is occupied then this looks at the next cell along and continues to do this until it finds an empty cell.  The code for this can be seen below where \texttt{pos} is the value calculated by the hash function.  As can be seen the cell \texttt{state} is used to detect whether a cell is occupied and \texttt{insertCell} is used to perform an actual insertion.  A for loop iterates over the sequence of cells looking for the next empty cell in the case of collisions.

Insert works by calculating a the hash of the string using my hash function.  It then looks at this cell in the hash table.  If the cell is empty then it places the value at that point in the cell -- if the cell is occupied then it uses the collision resolution strategy to find an alternative position for the value.

Find works in much the same way, but in this case it compares the cell's value.  If it matches the value of the string then it returns true.  Otherwise it uses the collision resolution strategy to find the next alternative cell and compares the value of that cell.  It continues in this fashion until an empty cell is encountered at which point it returns false.


In insert:
\ifc
\begin{verbatim}
	for(int j = 1; j < N; j++){
            int look = (pos + j) % N;
            set->collision ++;
            if(set->cells[look].state != in_use){
              insertCell(value, set, look);
              break;
            }
\end{verbatim}
\fi

\ifjava
\begin{verbatim}
if(isLinearProbing(mode)) // linear probing: A[(i+j)mod N]
   {
       for(int j = 1; j < N; j++){
            int look = (pos + j) % N;
            collision ++;
            if(cells[look].state != state.in_use){
                 insertCell(value, look);
                 break;
            }
      }
}
\end{verbatim}
\fi

\ifpython
\begin{verbatim}
if self.isLinearProbing(self.mode): # linear probing: A[(i + j)mod N]
    for j in range(1, N):
         look = (pos + j) % N
         self.collision = self.collision + 1
         if (self.cells[look].state != state.in_use):
                 self.insertCell(value, look)
                 break
         if (self.cells[look].element == value):
                 break
\end{verbatim}
\fi



In find:
\ifc
\begin{verbatim}
for(int j = 1; j < N; j++){
          int look = (pos + j) % N;
          if(set->cells[look].state == empty)
            return false;
          if(set->cells[look].state == in_use
             && (strcmp(set->cells[look].element, value) == 0)){
            return true;
          }
        }
\end{verbatim}
\fi

\ifjava
\begin{verbatim}
for(int j = 1; j < N; j++){
            int look = (pos + j) % N;
            if(cells[look] == null)
                    return false;
            if(cells[look].compareTo(value) == 0) {
                     return true;
             }
}
\end{verbatim}
\fi


\ifpython
\begin{verbatim}
for j in range(1, N):
                look = (pos + j) % N
                if (self.cells[look] == None):
                    return False
                if (self.cells[look] != None and self.cells[look] == value):
                    return True
\end{verbatim}
\fi

The advantage of linear probing is that it is simple to implement and effective if the table is a good size (relative to the input) and the hash function creates a good spread of values given the input.  It's disadvantage is that as the table fills up (particularly if the hash function isn't generating a good spread of values), \emph{clusters} will form which cause large blocks of the hash table to be fully occupied and then the function iterates the whole length of the block to find a space -- this can cause insertion and find to have linear complexity.

{\bf Marking Feedback:}  Most of the marks here got lost for silly mistakes - e.g., not actually describing the implementation.  A few people failed to either take the modulus of \texttt{pos + j} or alternatively to make sure that the value wrapped around to zero when the hash table size was reached.  This then had implications for the load factor and whether a position for an element once the table was full.  This often showed up in the tests with smaller initial dictionary sizes.

\begin{quote}
Did you implement rehashing? If so briefly explain your implementation including explaining when you choose to rehash.  (No more than 200 words, you may include code snippets which do not count towards the word count)
\end{quote}

Yes.  I rehash when the table is over half full ($num\_entries \geq N/2$) which is the point when a lot of collisions will be occurring.  The process works by calculating a new size for the hash table by doubling the old size and then finding the next prime.  A copy is made of the old hash table to store the values.  Then I reinitialise the existing hash table to the new size followed by iterating over the keys in the copy of the old hash table and inserting any values found into the new hash table.

{\bf Marking Feedback:}  A lot of people lost half a mark here for failing to find the next prime number after doubling the size of the hash table. 

\section{Part B}

\begin{quote}
How do your find and insert functions work using binary trees? How have you implemented them? Include your code for insert and find and relate your explanation to this code.  The code may be tidied up a bit for presentation reasons, but should include all important details mentioned in your explanation.  (No more than 200 words, not including the code snippet(s))
\end{quote}

\medskip 

Binary trees work by comparing strings using alphabetical order and extending a tree in which the trees are stored.  Each node in the tree stores a string as a value.  When a new value is to be inserted it is compared with the value at the current node, if it is less than the current value then it is inserted into the left sub-tree (line 3 below) and if it is greater than the current value then it is inserted into the right subtree (line \ifc 4 \fi \ifjava 5 \fi \ifpython 5 \fi below).  If there are no sub-trees then a new node is created at the leaf with empty sub-trees (the else statement below).

\ifc
Insert:
\begin{verbatim}
  if(tree){
    int comp = compare(value,tree->value);
    if(comp < 0){ tree->left = insert(value,tree->left); }
    else if(comp > 0){ tree->right = insert(value,tree->right); }
    // if comp == 0 then do nothing as it's a duplicate
  }
  else{
    tree = malloc(sizeof(struct bstree));
    check(tree);
    tree->left  = NULL;
    tree->right = NULL;
    tree->value= strdup(value);
  }
  return tree;
\end{verbatim}
\fi

\ifjava
\begin{verbatim}
if(tree()){
            int comp = value.compareTo(this.value);
            if (comp < 0) {
                this.left.insert(value);
            } else if (comp > 0) {
                this.right.insert(value);
            }
            // if comp == 0 then do nothing as it's a duplicate
        }
        else{
            left = new bstree(this.config);
            right = new bstree(this.config);
            this.value = value;
        }
\end{verbatim}
\fi

\ifpython
\begin{verbatim}
def insert(self, value):
        if (self.tree()):
            if(value < self.value):
                self.left.insert(value)
            elif (value > self.value):
                self.right.insert(value)
        else:
            self.left = bstree();
            self.right = bstree();
            self.value = value;
\end{verbatim}
\fi



Find:
\ifc
\begin{verbatim}
if(tree){
    int comp = compare(value,tree->value);
    if(comp==0){ return true; }
    else if(comp < 0){ return find(value,tree->left); }
    else { return find(value,tree->right); }
  }
  // if tree is NULL then it contains no values
  return false;
\end{verbatim}
\fi

\ifjava
\begin{verbatim}
if(tree()){
            int comp = value.compareTo(this.value);
            if (comp == 0) { return true;}
            else if (comp < 0) {return left.find(value);}
            else return right.find(value);
        }
        // if tree is NULL then it contains no values
        return false;
\end{verbatim}
\fi

\ifpython
\begin{verbatim}
def find(self, value):
        if self.tree():
            if (value == self.value):
                return True
            elif (value < self.value):
                return self.left.find(value)
            else:
                return self.right.find(value)
        return False
\end{verbatim}
\fi

Find works in much the same way.  The tree is explored moving left if the word to be found comes earlier in the alphabet than the search word and right otherwise.  If a leaf node is reached and the word has not been found then false is returned.

{\bf Marking Feedback:}  Most people got full marks here.  Where marks got lost it was generally for silly mistakes - e.g., not actually describing the implementation.  


\section{Part C}

\subsection{Experiment 1}

\paragraph{Hypothesis} The complexity of inserting a single value into my implementation of a hash set is $O(1)$. 

Theoretically, the best case for hash set insertion is $O(1)$ since a) the cost of inserting into a known point in a hash table is constant, b) calculation of the hash value should be constant and, c) on average, we don't expect collisions to occur.  We would expect average case behaviour to be similar provided we have a decent hash function and, in particular, if we have a big enough hash set to avoid rehashing.  

{\bf Marking Feedback:}. A lot of people lost half a mark here for stating that the best case was $O(1)$ without stating that this was because the cost of inserting into a known point in a hash table was constant.  People also lost marks for failing to state what was being counted as the input size (even though the complexity is constant time, we want to know with respect to what).

\paragraph{Experimental Design} 

To test the hypothesis we investigated the performance of the spell-checking program on random input dictionaries.  Our independent variable was the size of the dictionary which is the same as the number of items to be inserted.  Our dependent variable was the time taken to perform the insertion.

Random dictionaries of sizes 10K, 20K, 30K, 40K, 50K, 60K, 70K, 80K, 90K and 100K were generated.  Five dictionaries of each size were created.

An empty query file was used to avoid the time taken for queries impacting on running time.  The size of the initial hashset was set to 200K to avoid rehashing.  

The spell checking program was then run once on each dictionary with the empty query file.  The time for the program to execute was measured using the UNIX \texttt{time} command summing the \texttt{user} and \texttt{sys} values output by the command.  

The process of generating dictionaries and computing the time was automated using the shell scripts shown in Appendix~\ref{app:shell_scripts_exp1}.

The results were then plotted using \texttt{gnuplot} and \texttt{gnuplot}'s {\bf fit} functionality was used to fit a line $f(x) = m*x + q$  calculating values for $m$ and $q$ in the process.  $m$ should give the average insert time per insert if the complexity of insertion is $O(1)$.  


\paragraph{Results} The results are shown in Figure~\ref{fig:hash_insert}.
\ifc
\begin{figure}[htb]
\begin{center}
\begin{gnuplot}[terminal=jpeg, terminaloptions={size 400,300 font "Arial,9"}]
max(x, y) = (x > y ? x:y)
min(x, y) = (x < y ? x:y)
set title "Time Taken to insert Words into a Dictionary using a Hash Set data structure"
set ylabel "Time in Milliseconds"
set xlabel "Size of Dictionary"
set xrange [0:105000]

f(x) = m *x + q

fit f(x) 'data/exp1_c_data.dat' using  1:(1000*$2) via m, q

mq_value = sprintf("Parameter value\nm = %f\nq = %f", m, q)
set label 1 at 25000,140 mq_value

plot "data/exp1_c_data.dat" using 1:(1000*$2) t "time", f(x) ls 2 t 'Best fit line'
set out
\end{gnuplot}
\end{center}
\caption{Time taken to insert $n$ words into a dictionary implemented using a hash set, with best fit line $f(x) = mx + q$ shown and values for the parameters $m$ and $q$}
\label{fig:hash_insert}
\end{figure}
\fi

\ifjava
\begin{figure}[htb]
\begin{center}
\begin{gnuplot}[terminal=jpeg, terminaloptions={size 400,300 font "Arial,9"}]
max(x, y) = (x > y ? x:y)
min(x, y) = (x < y ? x:y)
set title "Time Taken to insert all Words into the Dictionary using a Hash Set Data structure"
set ylabel "Time in Milliseconds"
set xlabel "Size of Dictionary"
set xrange [0:105000]

f(x) = m*x + q

fit f(x) 'data/exp1_java_data.dat' using  1:(1000*$2 ) via m, q

mq_value = sprintf("Parameter value\nm = %f\nq = % f", m, q)
set label 1 at 25000,400 mq_value

plot "data/exp1_java_data.dat" using 1:(1000*$2) t "time", f(x) ls 2 t 'Best fit line'
set out
\end{gnuplot}
\end{center}
\caption{Time taken to insert $n$ words into a dictionary implemented using a hash set, with best fit line $f(x) = mx + q$ shown and values for the parameters $m$ and $q$}
\label{fig:hash_insert}
\end{figure}
\fi


\ifpython
\begin{figure}[htb]
\begin{center}
\begin{gnuplot}[terminal=jpeg, terminaloptions={size 400,300 font "Arial,9"}]
max(x, y) = (x > y ? x:y)
min(x, y) = (x < y ? x:y)
set title "Time Taken to insert all Words into the Dictionary"
set ylabel "Time in Milliseconds"
set xlabel "Size of Dictionary"
set xrange [0:105000]

f(x) = m*x + q

fit f(x) 'data/exp1_python_data.dat' using  1:(1000*$2 ) via m, q

mq_value = sprintf("Parameter value\nm = %f\nq = % f", m, q)
set label 1 at 2500,1000 mq_value

plot "data/exp1_python_data.dat" using 1:(1000*$2) t "time", f(x) ls 2 t 'Best fit line'
set out
\end{gnuplot}
\end{center}
\caption{Time taken to insert words into a dictionary implemented using a hash set, with best fit line $f(x) = mx + q$ shown and values for the parameters $m$ and $q$}
\label{fig:hash_insert}
\end{figure}
\fi

As can be seen from the graph the the line $f(x) = mx + 1$ has a good fit and the value of \ifc 0.0017 \fi \ifjava 0.01\fi \ifpython 0.01 \fi ms has been computed for $m$ and \ifc $-6.9$ \fi \ifjava 130 \fi \ifpython 8 \fi computed for $q$.  This confirms the hypothesis about the behaviour of the algorithm and allows us to predict the time taken to insert an item into the dictionary will be about \ifc 0.0016 \fi \ifjava 0.01 \fi \ifpython 0.01 \fi ms.  \ifc It should be noted that we have a negative value for $q$ which is clearly incorrect (small dictionaries will not be inserted in negative time).  This is probably the result of some noise in the data -- more runs of the program, possibly extending to larger dictionaries might have resolved this. \fi

{\bf Marking Feedback:} A lot of people got confused somewhere around results and discussion with the difference between the total time to insert $n$ items into a hash set (which we would expect to be $O(n)$ since each insert is constant) with the average time per insert (expected to be $O(1)$).   Exactly where marks started getting lost depended upon where this first became obvious -- having something linear showing in the graph was fine so long as it was clear this was total time or the text flagged this up, but then the discussion sometimes went on to say this disproved the hypothesis.


\subsection{Experiment 2}

\paragraph{Hypothesis}  The complexity of inserting a single value into my implementation of a binary tree is $O(log~n)$ where $n$ is the size of the dictionary. 

Theoretically, best case behaviour for binary trees is $O(log~n)$ when the tree is equally balanced, in this case $log~n$ is the height of the tree representing the length of the branch to be traversed in order to perform an insert.  We would expect average case behaviour to be similar to best case behaviour.  

\paragraph{Experimental Design} 

To test the hypothesis we investigated the performance of the spell-checking program on random input dictionaries.  Our independent variable was the size of the dictionary which is the same as the number of items to be inserted.  Our dependent variable was the time taken to perform the insertion.

Five random dictionaries of sizes \ifjava 1K, 2K, 3K, 4K, 5K, 6K, 7K, 8K, 9K,\fi \ifpython 1K, 2K, 3K, 4K, 5K, 6K, 7K, 8K, 9K,\fi  10K, 20K, 30K, 40K, 50K, 60K, 70K, 80K, 90K, and 100K  were generated.  

An empty query file was used to avoid the time taken for queries impacting on running time. 

The spell checking program was then run once on each dictionary with the empty query file.  The time for the program to execute was measured using the UNIX \texttt{time} command summing the \texttt{user} and \texttt{sys} values output by the command.  \ifjava The height of the resulting binary tree was computed by the program and also recorded. \fi

% The results were then plotted using \texttt{gnuplot}.   \ifjava There appeared to be a fixed overhead time for the program to execute that was independent of the size of the dictionary.  However this proved difficult to calculate from the data available, in part because there was a spread of running times for each size of dictionary.  Figure~\ref{fig:exp2_fixed} illustrates this spread of values.  The difficulty of calculating the overhead, meant it proved impossible to calculate time per insert accurately for a given dictionary size.  Instead, therefore, we broke the theoretical complexity down -- the theoretical complexity states that  the cost of insertion is $O(log~n)$ because it is linear in the height of the binary tree, and the height of the binary tree should be roughly $log n$ where $n$ is the size of the dictionary.  Therefore we computed time taken against the height of the tree (Figure~\ref{fig:exp2_time_height}) and the height of the tree against the size of the dictionary (Figure~\ref{fig:exp2_height_size}). \fi

For our calculations in Section~\ref{sec:conclusion} we will need to know the time taken to insert $n$ items into the dictionary.  This is easily plotted from the data collected and shown in Figure~\ref{fig:exp2_fixed} with the line  $f(x) = x \times m \times log~x + q$ shown.

\ifc In order to find the time per single insert, to check the hypothesis about the implementation, the time calculated to insert $n$ words into the dictionary was divided by the size of the dictionary ($n$). This is shown in Figure~\ref{fig:exp2}.\fi

\ifjava Figure~\ref{fig:exp2_fixed} shows there is a significant fixed overhead for any run of the program ($q = 100ms$). In order to compute time per single insert, needed to check the hypothesis about the implementation, we needed to plot $\frac{x - q}{n}$ where $x$ is the time taken to insert $n$ items into the dictionary against dictionary size and then attempt to fit the line $f(x) = m \times log~x' + q'$ to this.  This is shown in Figure~\ref{fig:exp2}\fi

\ifjava Figure~\ref{fig:exp2_fixed} shows there is a  fixed overhead for any run of the program ($q = 11ms$). In order to compute time per single insert, needed to check the hypothesis about the implementation, we needed to plot $\frac{x - q}{n}$ where $x$ is the time taken to insert $n$ items into the dictionary against dictionary size and then attempt to fit the line $f(x) = m \times log~x' + q'$ to this.  This is shown in Figure~\ref{fig:exp2}\fi

The process of generating dictionaries and computing the time was automated using the shell scripts shown in Appendix~\ref{app:shell_scripts_exp2}.

\paragraph{Results} 
Figure~\ref{fig:exp2_fixed} presents results of time to insert $n$ items into a random dictionary.  The line $f(x) = x m log~x + q$ has been calculated and the values for $m$ and $q$ are shown. \ifc
\begin{figure}[htb]
\begin{center}
\begin{gnuplot}[terminal=jpeg, terminaloptions={size 400,300 font "Arial,9"}]
max(x, y) = (x > y ? x:y)
min(x, y) = (x < y ? x:y)
set title "Time Taken to insert Words into a Dictionary using a Binary Tree"
set ylabel "Time in Milliseconds"
set xlabel "Size of Dictionary"
set xtics format "%.0s%c"
set xrange [0:90500]

logb(x, base) = log(x)/log(base)
f1(x) = m * logb(x, 2)  * x + q

fit f1(x) 'data/exp2_c_data.dat' using  1:((1000*$2)) via m, q

mq1_value = sprintf("Parameter value\nm = %f\nq = %f", m, q)
set label 1 at 8000,100 mq1_value

plot "data/exp2_c_data.dat" using 1:((1000*$2)) t "average case",  f1(x) ls 2 t 'Best fit line average case'
set out
\end{gnuplot}
\end{center}
\caption{Time taken to insert $n$ words into a dictionary implemented using a binary tree, with best fit line $f(x) = x \times m \times log~x + q$ shown and values for the parameters $m$ and $q$}
\label{fig:exp2_fixed}
\end{figure}

\begin{figure}[htb]
\begin{center}
\begin{gnuplot}[terminal=jpeg, terminaloptions={size 400,300 font "Arial,9"}]
max(x, y) = (x > y ? x:y)
min(x, y) = (x < y ? x:y)
set title "Average Time Taken to insert a Word into a Dictionary using a Binary Tree data structure"
set ylabel "Time in Milliseconds"
set xlabel "Size of Dictionary"
set xtics format "%.0s%c"
set xrange [0:105000]

logb(x, base) = log(x)/log(base)
f1(x) = m * logb(x, 2)  + q

fit f1(x) 'data/exp2_c_data.dat' using  1:(1000*$2/$1) via m, q

mq1_value = sprintf("Parameter value\nm = %f\nq = %f", m, q)
set label 1 at 50000,0.0012 mq1_value

plot "data/exp2_c_data.dat" using 1:(1000*$2/$1) t "average case",  f1(x) ls 2 t 'Best fit line average case'
set out
\end{gnuplot}
\end{center}
\caption{Average time taken to insert a single word into a dictionary implemented using a binary tree, with best fit line $f(x) = m log~x + q$ shown and values for the parameters $m$ and $q$}
\label{fig:exp2}
\end{figure}
\fi

\ifjava
\begin{figure}[htb]
\begin{center}
\begin{gnuplot}[terminal=jpeg, terminaloptions={size 400,300 font "Arial,9"}]
max(x, y) = (x > y ? x:y)
min(x, y) = (x < y ? x:y)
set title "Time Taken to insert all Words into the Dictionary using a Binary Tree data structure"
set ylabel "Time in Milliseconds"
set xlabel "Size of Dictionary"
set xtics format "%.0s%c"
set xrange [0:105000]

q=300
logb(x, base) = log(x)/log(base)
f1(x) = m * logb(x, 2)  * x + q

fit f1(x) 'data/exp2_java_data.dat' using  1:((1000*$2)) via m, q

mq1_value = sprintf("Parameter value\nm = %f\nq = %f", m, q)
set label 1 at 15000,250 mq1_value

plot "data/exp2_java_data.dat" using 1:((1000*$2)) t "average case",  f1(x) ls 2 t 'Best fit line average case'
set out

\end{gnuplot}
\end{center}
\caption{Time taken to insert $n$ words into a dictionary implemented using a binary tree, with best fit line $f(x) = x \times m \times log~x + q$ shown and values for the parameters $m$ and $q$}
\label{fig:exp2_fixed}
\end{figure}

\begin{figure}[htb]
\begin{center}
\begin{gnuplot}[terminal=jpeg, terminaloptions={size 400,300 font "Arial,9"}]
max(x, y) = (x > y ? x:y)
min(x, y) = (x < y ? x:y)
set title "Average Time Taken to insert a Word into a Dictionary using a Binary Tree data structure"
set ylabel "Time in Milliseconds"
set xlabel "Size of Dictionary"
set xtics format "%.0s%c"
set xrange [0:105000]

logb(x, base) = log(x)/log(base)
f1(x) = m * logb(x, 2)  + q

fit f1(x) 'data/exp2_java_data.dat' using  1:(1000*$2 - 110/$1) via m, q

mq1_value = sprintf("Parameter value\nm = %f\nq = %f", m, q)
set label 1 at 50000,250 mq1_value

plot "data/exp2_java_data.dat" using 1:(1000*$2 - 110/$1) t "average case",  f1(x) ls 2 t 'Best fit line average case'
set out
\end{gnuplot}
\end{center}
\caption{Average time taken to insert a single word into a dictionary implemented using a binary tree, with best fit line $f(x) = m log~x + q$ shown and values for the parameters $m$ and $q$}
\label{fig:exp2}
\end{figure}

\begin{figure}[htb]
\begin{center}
\begin{gnuplot}[terminal=jpeg, terminaloptions={size 400,300 font "Arial,9"}]
max(x, y) = (x > y ? x:y)
min(x, y) = (x < y ? x:y)
set title "Height of  Binary Tree data structure given size of Input"
set ylabel "Height of Tree"
set xlabel "Size of Input"
set xtics format "%.0s%c"

logb(x, base) = log(x)/log(base)
f1(x) = m * logb(x, 2)  + q
fit f1(x) 'data/exp2_java_data.dat' using  1:3 via m, q

mq1_value = sprintf("Parameter value\nm = %f\nq = %f", m, q)
set label 1 at 50000,20 mq1_value

plot "data/exp2_java_data.dat" using 1:3 t 'Height of Tree', f1(x) ls 2 t 'Best Fit Line'
set out
\end{gnuplot}
\end{center}
\caption{Average time taken to insert a single word into a dictionary implemented using a binary tree, with best fit line $f(x) = m log~x + q$ shown and values for the parameters $m$ and $q$}
\label{fig:exp2_height}
\end{figure}

\begin{figure}[htb]
\begin{center}
\begin{gnuplot}[terminal=jpeg, terminaloptions={size 400,300 font "Arial,9"}]
max(x, y) = (x > y ? x:y)
min(x, y) = (x < y ? x:y)
set title "Average Time Taken to insert a Word into a Dictionary using a Binary Tree data structure."
set ylabel "Time in Milliseconds"
set xlabel "Height of Tree"
set xtics format "%.0s%c"

f1(x) = m * x  + q
fit f1(x) 'data/exp2_java_data.dat' using  3:(1000*$2 - 110/$1) via m, q

mq1_value = sprintf("Parameter value\nm = %f\nq = %f", m, q)
set label 1 at 20,250 mq1_value

plot "data/exp2_java_data.dat" using 3:(1000*$2 - 110/$1) t 'time', f1(x) ls 2 t 'Best Fit Line'
set out
\end{gnuplot}
\end{center}
\caption{Average time taken to insert a single word into a dictionary implemented using a binary tree, with best fit line $f(x) = m log~x + q$ shown and values for the parameters $m$ and $q$}
\label{fig:exp2_time_height}
\end{figure}


% logb(x, base) = log(x)/log(base)
% f1(x) = m * x + q

% fit f1(x) 'data/exp2_java_data.dat' using  1:3 via m, q

% mq1_value = sprintf("Parameter value\nm = %f\nq = %f", m, q)
% set label 1 at 50000,0.0012 mq1_value

%f1(x) ls 2 t 'Best fit line average case'

\fi

\ifpython
\begin{figure}[htb]
\begin{center}
\begin{gnuplot}[terminal=jpeg, terminaloptions={size 400,300 font "Arial,9"}]
max(x, y) = (x > y ? x:y)
min(x, y) = (x < y ? x:y)
set title "Time Taken to insert Words into the Dictionary using a Binary Tree data structure"
set ylabel "Time in Milliseconds"
set xlabel "Size of Dictionary"
set xtics format "%.0s%c"
set xrange [0:105000]

logb(x, base) = log(x)/log(base)
f1(x) = m * logb(x, 2) * x + q

fit f1(x) 'data/exp2_python_data.dat' using  1:((1000*$2 )) via m, q

mq1_value = sprintf("Parameter value\nm = %f\nq = %f", m, q)
set label 1 at 6000,700 mq1_value

plot "data/exp2_python_data.dat" using 1:((1000*$2)) t "average case",  f1(x) ls 2 t 'Best fit line average case'
set out
\end{gnuplot}
\end{center}
\caption{Time taken to insert all words into a random dictionary, with best fit line $f(x) = x m log~x + q$ shown and values for the parameters $m$ and $q$}
\label{fig:exp2_fixed}
\end{figure}


\begin{figure}[htb]
\begin{center}
\begin{gnuplot}[terminal=jpeg, terminaloptions={size 400,300 font "Arial,9"}]
max(x, y) = (x > y ? x:y)
min(x, y) = (x < y ? x:y)
set title "Average Time Taken to insert a Word into the Dictionary using a Binary Tree data structure"
set ylabel "Time in Milliseconds"
set xlabel "Size of Dictionary"
set xtics format "%.0s%c"
set xrange [0:105000]

logb(x, base) = log(x)/log(base)
f1(x) = m * logb(x, 2)  + q

fit f1(x) 'data/exp2_python_data.dat' using  1:((1000*$2 - 11)/$1) via m, q

mq1_value = sprintf("Parameter value\nm = %f\nq = %f", m, q)
set label 1 at 60000,0.014 mq1_value

plot "data/exp2_python_data.dat" using 1:((1000*$2 - 11)/$1) t "average case",  f1(x) ls 2 t 'Best fit line average case'
set out
\end{gnuplot}
\end{center}
\caption{Time taken to insert words into a random dictionary, with best fit line $f(x) = m log~x + q$ shown and values for the parameters $m$ and $q$}
\label{fig:exp2}
\end{figure}

\fi

Figure~\ref{fig:exp2} presents the results for time to insert a single entry  a random dictionary.  The line $f(x) = m log~x + q$ has been calculated and the values for $m$ and $q$ are shown. \ifc As can be seen from the graph the the line $f(x) = m log~x + q$ has a reasonable fit in the average case and the value of 0.000252ms has been computed for $m$ and -0.002ms has been computed for $q$.  This confirms the hypothesis about the behaviour of the algorithm in the average case  is $O(log~n)$ -- though the values for the best fit line are obviously not quite right since $q$ should not be negative\footnote{For an 8 hour lab it is fine to leave this here as an observation that the best fit line needs some work.}. \fi  
\ifjava As can be seen from the graph the the line $f(x) = m log~x + q$ has an OK fit in the average case (particularly if the larger input files are ignored) and the value of 22ms has been computed for $m$ and -145ms has been computed for $q$.  This suggests the hypothesis about the behaviour of the algorithm in the average case  is $O(log~n)$ -- though a better fit would have been preferable -- in particular without a negative value for $q$\footnote{For an 8 hour lab it is fine to leave this here as an observation that the best fit line needs some work but, as you can see, I did a little more work to investigate this further.}.  

\fbox{\parbox{0.9\textwidth}{\textbf{What did I do about the bad results here?:} If you compare these Java answers with the C answers, you will see that the smaller input dictionaries (sizes under 10K) weren't used in the C case.  When I observed that I had something that looked more like a straight line in the average time, I hypothesised that I was looking at the part of the data where the log curve is flattening off, so I wanted to see what was happening at smaller dictionary sizes.  It's these results that are shown and they reassured me that my program wasn't completely wrong. 
}}

Theoretically the time taken in the average case is $O(log~n)$ because the height of the resulting tree should be $O(log~n)$ and the time should relate to the height of the tree.  In order to verify that the algorithm was working correctly I adapted the program to print out the height of the tree and Figure~\ref{fig:exp2_height} plots the height of the tree against the input.  With a best fit line of $f(x) = m log~x + q$ shown.  This is a much better fit than the line in Figure~\ref{fig:exp2} although it can be seen that there is a lot of variation in the tree height making it hard to distinguish the larger values.  Figure~\ref{fig:exp2_time_height} shows the time taken for the program to run against the height of the tree with the line $f(x) = mx + q$.  As can be seen this has a good fit for most of the values, but there are some anomalously slow results for some of the larger trees -- these probably correspond to the slower results observed for the larger dictionaries.  If more time was available, further investigation should re-run these results, possibly using a different timing mechanisms and a profiling tool, particularly for larger dictionaries to see if the anomalous slow down was caused garbage collection within the Java virtual machine or by factors other than the running time of the program itself.
\fi
 \ifpython As can be seen from the graph the the line $f(x) = m log~x + q$ has a reasonable fit in the average case and the value of 0.0005ms has been computed for $m$ and 11ms has been computed for $q$.  This confirms the hypothesis about the behaviour of the algorithm in the average case  is $O(log~n)$ though it should be noted there is one outlier for the smallest dictionary (see in Figure~\ref{fig:exp2}).
 
 \fbox{\parbox{0.9\textwidth}{\textbf{A note about what I did here:} If you compare these Python answers with the C answers, you will see that the smaller input dictionaries (sizes under 10K) weren't used in the C case.  When I observed that I had something that looked more like a straight line in the average time, I hypothesised that I was looking at the part of the data where the log curve is flattening off, so I wanted to see what was happening at smaller dictionary sizes.  It's these results that are shown and they reassured me that my program wasn't completely wrong. 
}}

 \fi %-- though the values for the best fit line are obviously not quite right since $q$ should not be negative\footnote{For an 8 hour lab it is fine to leave this here as an observation that the best fit line needs some work.}. \fi  

{\bf Marking Feedback:} A lot of people got confused somewhere around results and discussion with the difference between the total time to insert $n$ items into a binary tree (which we would expect to be $O(n~log~n)$ since each insert is $O(log~n)$) with the average time per insert.   Exactly where marks started getting lost depended upon where this first became obvious -- having something linear-ish showing in the graph was fine so long as it was clear this was total time or the text flagged this up, but then the discussion sometimes went on to say this disproved the hypothesis.  A number of people also lost marks by over-stating how good a fit their line was to the data.


\section{Conclusion}
\label{sec:conclusion}
Suppose we are considering the time take to build a dictionary?  The time taken to create a dictionary of $x$ items as a hash set is $m_h x + q_h$ and the time taken to create a dictionary of $x$ items as a binary tree is $m_b x log~x + q_b$.  The point where a hash set becomes preferable to a binary tree (in the average case) is therefore the point where 
$$m_h x + q_h < m_b x log~x + q_b$$
Assuming $q_h$ and $q_b$ are negligible this is the point where $m_h < m_b log~x$ which is where $\frac{m_h}{m_b} < log~x$.  In our case this is where \ifc $log~x > \frac{0.0017}{0.000116} \approx 15$ \fi 
\ifjava $log~x > \frac{0.01}{0.000092} \approx 108.6$ \fi. 
\ifpython  $log~x > \frac{0.013}{0.0004} \approx 32.5$  \fi 
So where \ifc $x > 2^{15}$ \fi
 \ifjava $x > 2^{108.6}$ \fi \ifpython $x > 2^{32.5}$ \fi which is somewhere around a dictionary size of \ifc 32k. \fi   \ifjava $5 \times 10^{32}$ which is large enough that for many dictionaries the use of a binary tree is preferable to a hash set. \fi  \ifc This is borne out by looking at Figure~\ref{fig:exp2} where we see that the time per insert into a binary tree exceeds 0.0017 (the time per insert in the hash set) at around this point. \fi     \ifpython 6,000,000,000 which is large enough that for many dictionaries the use of a binary tree is preferable to a hash set. \fi             

{\bf Marking Feedback:} People got marks here for either performing a calculation or superimposing their best fit lines and noting the crossover point.  A surprising number of people then got confused and said that hash sets should be used for smaller sets and binary trees for larger sets even though they had shown hash set complexity was constant while their binary tree complexity was $O(log~n)$.

\section{Extension}
For the extension I look at query time for hash sets.
 \subsection{Experiment 3}
\paragraph{Hypothesis} The complexity of querying a single value in my implementation of a hash set $O(1)$. 

Theoretically, the best case for hash set querying is $O(1)$ since the cost of looking up and inserting into a known point in a hash table is constant and, on average, we don't expect collisions to have occurred during insertion.  We would expect average case behaviour to be similar provided we have a decent hash function.  

\paragraph{Experimental Design} To test the hypothesis, random input dictionaries were produced. 5 unsorted dictionaries of sizes 10K, 20K, 30K, 40K and 50K  were generated.  For each random dictionary five random query files were generated of sizes 20K, 40K, 60K, 80K and 100K where each query file had 50\% of the words in the dictionary and 50\% of the words were not in the dictionary. The spell checking program was then run on each dictionary firstly with an empty input file  to calculate the time taken to create the dictionary and then with the query files.  The time for the program to execute was measured using the UNIX \texttt{time} command summing the \texttt{user} and \texttt{sys} values output by the command.  

 The results were then plotted using \texttt{gnuplot}.  In order to find the time per query the time calculated to make all the queries, less the time to create the dictionary, was divided by the size of the query file.  

The process of generating dictionaries and computing the time was automated using the shell scripts shown in Appendix~\ref{app:shell_scripts_exp3}.

\paragraph{Results}  The results of time per query are shown in Figure~\ref{fig:exp3} and the line $f(x) = m$ has been calculated and the value of $m$ is shown.

\ifc
\begin{figure}[htb]
\begin{center}
\begin{gnuplot}[terminal=jpeg, terminaloptions={size 400,300 font "Arial,9"}]
max(x, y) = (x > y ? x:y)
min(x, y) = (x < y ? x:y)
set title "Average Time per Query into a Hash Data Structure"
set ylabel "Time in Milliseconds"
set xlabel "Size of Dictionary"
set xrange [0:55000]
set yrange [0:0.005]

f(x) = m 

fit f(x) 'data/exp3_c_data.dat' using  1:(1000*($3 - $4)/$2) via m

mq_value = sprintf("Parameter value\nm = %f\n", m)
set label 1 at 25000,0.003 mq_value

plot "data/exp3_c_data.dat" using 1:(1000*($3 - $4)/$2) t "time", f(x) ls 2 t 'Best fit line'
set out
\end{gnuplot}
\end{center}
\caption{Time taken to query words against a dictionary implemented using a hash set, with best fit line $f(x) = m$ shown and value for the parameter $m$}
\label{fig:exp3}
\end{figure}
\fi

\ifjava
\begin{figure}[htb]
\begin{center}
\begin{gnuplot}[terminal=jpeg, terminaloptions={size 400,300 font "Arial,9"}]
max(x, y) = (x > y ? x:y)
min(x, y) = (x < y ? x:y)
set title "Average Time per Query"
set ylabel "Time in Milliseconds"
set xlabel "Size of Dictionary"
set xrange [0:55000]
set yrange [0:0.03]

f(x) = m 

fit f(x) 'data/exp3_java_data.dat' using  1:(1000*($3 - $4)/$2) via m

mq_value = sprintf("Parameter value\nm = %f\n", m)
set label 1 at 2500,0.025 mq_value

plot "data/exp3_java_data.dat" using 1:(1000*($3 - $4)/$2) t "time", f(x) ls 2 t 'Best fit line'
set out
\end{gnuplot}
\end{center}
\caption{Time taken to query words against a dictionary, with best fit line $f(x) = m$ shown and value for the parameters $m$}
\label{fig:exp3}
\end{figure}
\fi

\ifpython
\begin{figure}[htb]
\begin{center}
\begin{gnuplot}[terminal=jpeg, terminaloptions={size 400,300 font "Arial,9"}]
max(x, y) = (x > y ? x:y)
min(x, y) = (x < y ? x:y)
set title "Average Time per Query"
set ylabel "Time in Milliseconds"
set xlabel "Size of Dictionary"
set xrange [0:55000]
set yrange [0:0.02]

f(x) = m 

fit f(x) 'data/exp3_python_data.dat' using  1:(1000*($3 - $4)/$2) via m

mq_value = sprintf("Parameter value\nm = %f\n", m)
set label 1 at 2500,0.005 mq_value

plot "data/exp3_python_data.dat" using 1:(1000*($3 - $4)/$2) t "time", f(x) ls 2 t 'Best fit line'
set out
\end{gnuplot}
\end{center}
\caption{Time taken to query words against a dictionary, with best fit line $f(x) = m$ shown and value for the parameters $m$}
\label{fig:exp3}
\end{figure}
\fi

As can be seen from the graph the the line $f(x) = m$ has a good fit in the average case and the value of \ifc 0.0014 \fi \ifjava 0.008 \fi \ifpython 0.008 \fi ms has been computed for $m$.  This confirms the hypothesis about the behaviour of the algorithm in the average case.

{\bf Marking Feedback:} Note this is not a first class answer - a first class answer would have attempted something that needed a different approach to the experimental design or significant extra implementation.  However this answer would get Satisfactory.
 
 \section{Data Statement}

The raw data for experiment 1 can be seen in Appendix~\ref{app:data_exp1}.
The raw data for experiment 2 can be seen in Appendix~\ref{app:data_exp2}.
The raw data for experiment 2can be seen in Appendix~\ref{app:data_exp3}.

All the raw data is also available in the \texttt{data} directory of the \texttt{model\_solutions} folder on gitlab.

\newpage 
\appendix

\section{Raw Data for Experiment 1}
\label{app:data_exp1}

\ifc
\csvreader[
tabular={|l|l|},
table head={\hline \multicolumn{2}{|c|}{Data} \\ \hline Size & Time (s)\\\hline},
late after line=\\\hline,
no head]{data/exp1_c_data.csv}{}%
{\csvcoli & \csvcolii}
\fi

\ifjava
\csvreader[
tabular={|l|l|},
table head={\hline \multicolumn{2}{|c|}{Data} \\ \hline Size & Time (s)\\\hline},
late after line=\\\hline,
no head]{data/exp1_java_data.csv}{}%
{\csvcoli & \csvcolii}
\fi

\ifpython
\csvreader[
tabular={|l|l|},
table head={\hline \multicolumn{2}{|c|}{Data} \\ \hline Size & Time (s)\\\hline},
late after line=\\\hline,
no head]{data/exp1_python_data.csv}{}%
{\csvcoli & \csvcolii}
\fi

\section{Shell Scripts for Experiment 1}
\label{app:shell_scripts_exp1}
\subsection{Generating Dictionaries}
\verbatiminput{dict_query_generate_1.sh}


\subsection{Script for Computing Run Times}
\ifc
\verbatiminput{data_generate_exp1_c.sh}
\fi

\ifjava
\verbatiminput{data_generate_exp1_java.sh}
\fi

\ifpython
\verbatiminput{data_generate_exp1_python.sh}
\fi


\section{Raw Data for Experiment 2}
\label{app:data_exp2}
\ifc
\csvreader[
tabular={|l|l|},
table head={\hline \multicolumn{2}{|c|}{Data} \\ \hline Size & Time (s)\\\hline},
late after line=\\\hline,
no head]{data/exp2_c_data.csv}{}%
{\csvcoli & \csvcolii}
\fi

\ifjava
\csvreader[
tabular={|l|l|},
table head={\hline \multicolumn{2}{|c|}{Data} \\ \hline Size & Time (s)\\\hline},
late after line=\\\hline,
no head]{data/exp2_java_data.csv}{}%
{\csvcoli & \csvcolii}
\fi

\ifpython
\csvreader[
tabular={|l|l|},
table head={\hline \multicolumn{2}{|c|}{Data} \\ \hline Size & Time (s)\\\hline},
late after line=\\\hline,
no head]{data/exp2_python_data.csv}{}%
{\csvcoli & \csvcolii}
\fi 


\section{Shell Scripts for Experiment 2}
\label{app:shell_scripts_exp2}

\ifc
\verbatiminput{data_generate_exp2_c.sh}

\fi

\ifjava
\verbatiminput{data_generate_exp2_java_height.sh}

\fi

\ifpython
\verbatiminput{data_generate_exp2_python.sh}

\fi

\section{Raw Data for Experiment 3}
\label{app:data_exp3}

\ifc
\csvreader[
tabular={|l|l|l|l|},
table head={\hline \multicolumn{4}{|c|}{Data} \\ \hline Dictionary Size & Number of Queries & Total Time & Time to Generate Dictionary \\\hline},
late after line=\\\hline,
no head]{data/exp3_c_data.csv}{}%
{\csvcoli & \csvcolii & \csvcoliii & \csvcoliv}
\fi

\ifjava
\csvreader[
tabular={|l|l|l|l|},
table head={\hline \multicolumn{4}{|c|}{Data} \\ \hline Dictionary Size & Number of Queries & Total Time & Time to Generate Dictionary \\\hline},
late after line=\\\hline,
no head]{data/exp3_java_data.csv}{}%
{\csvcoli & \csvcolii & \csvcoliii & \csvcoliv}
\fi

\ifpython
\csvreader[
tabular={|l|l|l|l|},
table head={\hline \multicolumn{4}{|c|}{Data} \\ \hline Dictionary Size & Number of Queries & Total Time & Time to Generate Dictionary \\\hline},
late after line=\\\hline,
no head]{data/exp3_python_data.csv}{}%
{\csvcoli & \csvcolii & \csvcoliii & \csvcoliv}
\fi

\section{Shell Scripts for Experiment 3}
\label{app:shell_scripts_exp3}
\subsection{Generating Dictionaries}
\verbatiminput{dict_query_generate_4.sh}

\subsection{Script for Computing Run Times}
\ifc
\verbatiminput{data_generate_exp3_c.sh}
\fi

\ifjava
\verbatiminput{data_generate_exp3_java.sh}
\fi

\ifpython
\verbatiminput{data_generate_exp3_python.sh}
\fi



\end{document}